% M. S. Tsoeu (2011), University of Cape Town <mohohlo.tsoeu@uct.ac.za>

% This is a project report templace document created for EEE4022FS students at the University of Cape Town.
%
% This file should be is processed with ``pdflatex`` and might need a few modifications if a different processor is chosen.


\documentclass[a4paper,12pt]{report}

%Include packages you need to use here

\usepackage[top = 1in, bottom = 1in, left = 1in, right = 1in]{geometry}
\usepackage{graphicx}
\usepackage{fancyhdr}
\usepackage{amsmath, amsthm, amssymb}
\usepackage{lastpage}
\usepackage{subfigure}
\usepackage{lscape}
\usepackage{hyphenat}
\usepackage{setspace}
\usepackage{hyperref}
%\usepackage[inline]{enumitem}   

%\usepackage{longtable}
\usepackage{multirow}
\usepackage{amsmath}
\usepackage{listings}
\usepackage{color}
%\usepackage{subcaption}
%\usepackage[export]{wrapfigure}

\definecolor{codegreen}{rgb}{0,0.6,0}
\definecolor{codegray}{rgb}{0.5,0.5,0.5}
\definecolor{codepurple}{rgb}{0.58,0,0.82}
\definecolor{backcolour}{rgb}{0.95,0.95,0.92}

\lstdefinestyle{mystyle}{
	backgroundcolor=\color{backcolour},   
	commentstyle=\color{codegreen},
	keywordstyle=\color{magenta},
	numberstyle=\tiny\color{codegray},
	stringstyle=\color{codepurple},
	basicstyle=\footnotesize,
	breakatwhitespace=false,         
	breaklines=true,                 
	captionpos=b,                    
	keepspaces=true,                 
	numbers=left,                    
	numbersep=5pt,                  
	showspaces=false,                
	showstringspaces=false,
	showtabs=false,                  
	tabsize=2
}

\lstset{style=mystyle}

% Include page formatting here. 
\parskip = 6mm
\parindent = 0mm
\renewcommand{\headrulewidth}{0pt}
\rhead[]{\thesection}
\lhead[\thechapter]{}


\begin{document}
 
% This section formats the title page of the Report.
\thispagestyle{empty}
{\Huge \begin{center}
% Modify the line below to insert your title.
Gait sequence modelling and estimation
\hrule 
% Modify the line below to insert your subtitle.
{\Large using Hidden Markov Models}
\end{center}}

\vskip 5mm
\begin{center}
\- \- \- \- \- \- \- \- \- \-\includegraphics[scale = 0.3]{uctLogo.png}
\end{center}

\vskip 5mm
\begin{center}
Prepared by:\\
\textbf{Kouame Kouassi}		% Insert your name here
\end{center}

\vskip 10mm
\begin{center}
Prepared for:\\
\textbf{A. Prof. Fred Nicolls}\\ 		% Insert your supervisor's name here.
Dept. of Electrical and Electronics Engineering\\University of Cape Town
\end{center}


\vskip 10mm
\begin{center}
Submitted to the Department of Electrical Engineering at the University of Cape Town in partial
fulfilment of the academic requirements for a Bachelor of Science degree in \textbf{Electrical and Computer Engineering}
\end{center}

\vskip 10mm
\textbf{Key Words:} Hidden Markov Models, gait estimation, quadruped, IMU, Gaussian mixtures, dimensionality reduction


\vskip 5mm
\begin{center}{\bf \today}
\end{center}

\newpage
\thispagestyle{empty}
\mbox{}
\newpage

\onehalfspacing
\nohyphens{
\thispagestyle{empty}
\vskip 40mm


% Please leave the declaration as it is (Standard UCT declaration).
{\Large Declaration}\\
\hrule

\vskip 10mm
\begin{enumerate}
\item I know that plagiarism is wrong. Plagiarism is to use another's work and pretend that it is one's
own.
\item I have used the IEEE convention for citation and referencing. Each contribution to, and quotation in,
this report from the work(s) of other people has been attributed, and has been cited and
referenced.
\item This report is my own work.
\item I have not allowed, and will not allow, anyone to copy my work with the intention of passing it off
as their own work or part thereof.
\end{enumerate}
\vskip 10mm
Signature:\ldots\ldots\ldots\ldots\ldots\ldots\ldots\ldots\ldots 
\\Kouame H. Kouassi 		% Chante this line to your name.
\vskip10mm
Date:\ldots\ldots\ldots\ldots\ldots\ldots\ldots\ldots\ldots\ldots .


\fancyfoot[C]{\thepage}
\pagestyle{plain}
\newpage
\pagenumbering{roman}
{\Large Acknowledgments}\\
\hrule
God - For unfailing leading and supporting me throughout this project.

A. Prof Nicolls, my supervisor - for great supervision, advice, insight, open-door policy, and a workspace in the image processing laboratory.

Dr. Amir Patel - For the IMU dataset.

PEJ, my scholarship programme - For giving me the opportunity to study at UCT.

Family - For motivation and their remote support.

Friends - For their encouragement and understanding.

My colleagues in the lab: For a conducive work environment, passionate discussions and their interest in my progress.
\newpage

{\Large Abstract}\\
\hrule

% Place your abstract here.
Hidden Markov Models (HMMs) are doubly embedded stochastic processes with a rich underlying statistical structure. 
In the present work, gait sequence modelling and estimation was performed with HMMs using inertial measurement unit (IMU) data. More specifically, from IMU measurements, a dog's footfalls were correctly identified with up to 95\% precision. The continuous-valued IMU measurements were modelled with Gaussian mixtures. The originality of this project lies in the following three reasons. \\\\
Firstly, the similar work found in literature only uses 6-dimensional accelerometer and gyroscope measurements \cite{ches2012} \cite{cont2013}. The gait estimation was performed using 18-dimensional data. Thus, feature extraction and feature subset selection techniques were used for dimensionality reduction. The experiments demonstrated that up to 78\% performance increase is obtainable by using dimensionality reduction when the training data is not large.\\\\
Secondly, the available IMU dataset was very small in size. Two different methods were employed to overcome this limitation. \\In the first method,
the reverse gait measurements were appended to the original data to increase the training data. In addition, the IMU measurements of the dog running, walking, and trotting were aggregated. 
Subsequently, it was possible to build more than 90\% accurate HMM model with 18-dimensional IMU measurements with low variance error.\\\\
Thirdly, the HMM gait estimation model was designed by building separate models for the front and the back legs of the dog. It, therefore, follows that the gait estimation algorithm developed in this project is applicable to both quadrupeds and bipeds.\\\\
Finally, the designed HMM models were used to successfully perform motion type recognition. Overall, this undergraduate final year project may be deemed successful.
\newpage
\tableofcontents

\newpage
\listoffigures

\newpage
\listoftables

% Page formatting, to place section titles as headers of odd pages and Chapter titles as headers of even pages.
\newpage
\fancyhead[RE,LO]{}
\fancyhead[LE]{\leftmark}
\fancyhead[RO]{\rightmark}
\pagestyle{fancy}

\pagenumbering{arabic}

% THe files included below are .tex files containing the respective chapters these are already created in this package and you can add to or modify them.
\chapter{Introduction}

Human motion analysis is an important research topic motivated by medical, artistic or scientific purposes \cite{towa2009}. Among others reasons, it is investigated for human activity monitoring, fall detection, gesture recognition, balance control evaluation, and abnormal behaviour detection \cite{cont2013}. To carry out these studies, micro-electro-mechanical inertial measurement units (IMUs) are increasingly used to capture motion data. Primarily, because they are lightweight and cheap \cite{ches2012} \cite{towa2009} \cite{cont2013}.\\
Gait analysis requires the identification of latent gait states \cite{cont2013}. A gait state being a particular stance in a given gait activity.
For this reason, it is extensively used in biologically inspired robot design. More specifically, understanding animals gait mechanism is very invaluable in bio-robotics.\\
Although, human gait analysis has been successfully performed using IMU, the threshold-based algorithms, %%TODO: ref 30 in "COnt'
and the fuzzy logic methods %%TODO: ref 8 in "COnt' 
used, lack robustnes \cite{cont2013}. Thus, a more robust probabilistic technique is investigated in the present work. In this final year undergraduate project, we use Hidden Markov Model (HMM) to perform gait sequence modelling and estimation with IMU measurements of a moving dog.


\section{Background to the study}
Bio-inspired robotics identifies useful mechanisms from nature to inform real-world engineering systems. In this vein, the University of Cape Town (UCT) set out to investigate how a cheetah uses its tail for stability during high acceleration, quick turns and sudden braking for sophisticated robots design. \\
To achieve this purpose, inertial measurement unit (IMU) data for a dog running, walking and trotting was acquired and labelled with the corresponding footfalls.\\
This dataset can, therefore, be used to perform gait sequence modelling and analysis using Hidden Markov Models (HMM), given its relatively small size. 	%%TODO: add to gain more insight into a dog's gait to mimic a cheetah					

\section{Objectives of this study}
The objective of the present work as stated in its outline is to "design, implement, and test (Hidden Markov Models) for estimating the gait sequence
from IMU data, so that specific models can be formulated and their
parameters estimated and interrogated." Thus, four main sub-objectives are formulated.
\begin{itemize}
	\item Formulate HMM models to estimate the gait sequence of a dog from the labelled IMU measurements.
	\item Implement, test and evaluate the formulated models in order to estimate their parameters.
	\item Investigate how dimensionality reduction affects the performance of the designed models.
	\item Using the implemented models, perform useful gait pattern analysis such as motion recognition.
\end{itemize}

To achieve the above objectives the hypotheses have been formulated:
\begin{itemize}
	\item HMMs can successfully model gait sequence dynamics using IMU data, in the absence of huge training data.
	\item Data aggregation and mirroring techniques can be used to overcome training data size limitations.
	\item Dimensionality reduction can increase the performance of an HMM, in the absence of a large training set.
	\item HMMs can be used to successfully perform gait activity recognition.
\end{itemize}


\section{Scope and limitations}
The project was proposed with a very broad scope. In order to ensure completion of the project in the allocated time, its scope was therefore reduced to the most essential aspects.\\
Thus the project focues on the design, the implementation, and the performance evaluation of HMM models to identify the dog's footfalls from IMU measurements.\\
Moreover, particular attention was paid to the design of dimensionality reduction techniques and how they affect the performance of the HMM models when the training dataset is not large.\\
Furthermore, using the implemented models, motion type recognition was performed using the implemented HMM models. However, the capturing and calibration of the IMU data is not part of the scope of this project.

The main constraint of this project was the relatively small size of the available IMU data.
\section{Plan of development}
Here you tell the reader how your report has been organised and what is included in each
chapter.

\chapter{Literature Review}
\section{Gait sequence modelling and estimation}

\subsection{Quadrupede gait modelling}
\subsubsection{Periodicity}
\subsection{Quadrupede gait estimation}

\section{Hidden Markov Models}
Hidden Markov Models (HMMs) are doubly embedded stochastic processes with a rich underlying statistical structure. Introduced at the end of the 1960s by Baum and colleagues, they have become one of the prefered techniques in speech recognition after the implementation of Baker and Jelinek in 1970s. HMMs have been successfully applied to various other engineering problems in pattern recognition for classification and fraud detection purposes, amongst others \cite{tuto1989} \cite{tool2001} \cite{towa2009} \cite{twop2008}.\\
The type of HMM depends on the possible connections between the states. Thus, an HMM in which a state can transition to any other state is an ergodic. Other types such as the Left-Right model or Bakis do not allow all possible transitions between the states.
\subsection{HMM parameters specification}
An HMM is fully specified by the following parameters 
\begin{enumerate}
	\item N, the number of distinct states of the model. Together they form the set of individual states \(S = \{S_1, S_2, ..., S_N\}\).
	\item T, the number of observations. A sample observation sequence is denoted as \(O = \{O_1, O_2, ..., O_T\}\).
	\item \(Q = {q_t}\), the set of states with \(q_t\) denoting the current state at time instance, t such that \(q_t \epsilon S\) and \(t = 1, 2, ..., T\).
	\item K, the number of distinct observation symbols per state. 
	\item \(V = \{v_1, v_2, ..., v_K\}\), the feature set of K dimensions.
	\item \(A =  \{a_{ij} \}\), the state transition probabilities. \(a_{ij}\)  denotes probability of transitioning from state \(S_i\) to state \(S_j\).
	\item \(\Phi =   \{ \phi_{j}(k\}\), the probability distribution of observation symbols in state j.
	\item The initial state distribution, \(\pi = {\pi_i}\)
\end{enumerate}
For continuous HMM (CHMM), i.e, HMM with continuous-valued observations, \(\Phi\) consists in a probability distribution function. Many applications have succefully modelled such distributions with mixtures of Gaussian distributions \cite{cont2013} \cite{ches2012}.
As such, \(\phi\) is approximated by a weighted sum of M multivariate Gaussian distributions \(\eta\). For a given, observation sequence, \(\phi\)  and \(\eta\) are therefore given by equations \ref{eq:phi} and \ref{eq:eta},

\begin{align} 
	\phi(O_t) = \sum_{m=1}^M \beta_{jm} \eta(\mu_{jm}, \Sigma_{jm}, O_t), \label{eq:phi} \\
	\eta(\mu, \Sigma, O) = \frac{1}{\sqrt{(2\pi)^K|\Sigma|}}exp(-\frac{1}{2}(O-\mu)'\Sigma^{-1}(O-\mu) \label{eq:eta}
\end{align} 
\begin{align*}
	1 \leq j \leq N; 1 \leq m  \leq M; \beta_{jm} \geq 0; \sum_{m=1}^{M}\beta_{jm} = 1
\end{align*}
where \(\beta_{jm}\) is the mixture composition coefficient; \(\mu_{jm}\), \(\Sigma_{jm}\), respectively the mean vector and covariance matrix of state j; M is the number of mixture components and K is the dimensionality of O. In practice, the log-likelihood of \(\phi(O_t)\) is rather computed to avoid overflow when implemented on a machine. %%TODO: ref this fact

As a summary, the compact specification of a continous valued observation HMM is defined by \ref{eq:CHMM} and that of a discrete HMM in \ref{eq:DHMM}.
\begin{align} 
	CHMM = \lambda_C = (A, \beta_{jm}, \mu_{jm}, \Sigma_{jm}, \pi) \label{eq:CHMM} \\
	DHMM = \lambda_D = (A, b_j(k), \pi) \label{eq:DHMM}
\end{align}

\subsubsection{Basic assumptions of HMMs theory}
HMM theory is built on three basic assumptions listed below.
\begin{enumerate}
\label{con:markov} \item \textit{The Markov assumption}: HMM assumes that the probability of being in the current	at any instance of time t, is uniquely dependent on the previous state, at time, t + 1. More specifically, \(a_{ij} = P[q_t = S_j|q_{t+1}=S_i]\). This assumption makes it unsuitable for long-range correlation capturing applications.

\item \textit{The stationary assumption}: Furthermore, HMM state transition probabilities are assumed to be time-independent. Thus, the transition probabilities of two distinct time, \(t_1\) and \(t_2\) are identical, \(P[q_{t_1} = S_j|q_{t_1 -1} = S_i] = P[q{t_2}=S_i|q_{t_2-1} = S_i]\). HMMs can therefore effectively model mechanisms with stationary observations.

\item \textit{The output/observation independence assumption}: The current observation also known as emission symbol is statistically independent of the previous observations. It is "emitted" only by the current state, \(P[O|q_1, q_2, ..., q_T, \lambda]=\prod_{t=1}^{T}P[O_t|q_t, \lambda]\).
\end{enumerate} 
The three assumptions make an HMM model a relatively simple graphic modelling to be implemented. This simplicity naturally comes with some limitations in modelling more complex problems, which however, may be modelled with higher order HMMs. %%TODO: look for ref for this fact
Futhermore, the three assumptions are very similar to those of a Markov chain. This is because the stochastic process of an HMM pertaining to the hidden states can be reduced to a Markov chain. In fact, an HMM is an extension of a Markov Chain. The essential difference between the two is that, with the former, there is no a one-to-one mapping between the states and the observation symbols \cite{biol1998}.

\subsection{The three basics problems for HMM design}
In \cite{tuto1989}, Lawrence argued that an HMM design needs to answer three fondamental problems. They are the \textit{training problem}, the \textit{evaluation problem}, and the \textit{decoding problem}. Each problem and its solution is discussed in greater details next. 

\subsubsection{The evaluation problem}
The evaluation problem is about answering this question:
\textit{Given the observation sequence \(O = O_1O_2O_T\), and a model \(\lambda\), how do we efficiently compute \(P(O|\lambda\), the probability of the observation sequence?} \cite{tuto1989} \cite{ches2012} \cite{cont2013}.
The naive answer to this question is simply computing the \(P(O|\lambda)\) according to equation \ref{eq:P}:
\begin{align}
	P(O|\lambda) = \sum_{q_1}^{q_T}\pi_{q_1}b_{q_1}(O_1)a_{q_1q_2}b_{q_2}(O_2)...a_{q_{T-1}q_T}(O)) \label{eq:P}
\end{align}
This approach has two issues, it is not only, computationally too expensive because of the exponential complexity with respect to T, but also, intractable for very long sequence. In pactice, \(P(O|\lambda)\) is computed by an algorithm called \textit{forward-backward} procedure, which is a more efficient method.

\subsubsection{The decoding problem}
 The decoding problem can be reduced to this interrogation: \textit{Given the observation sequence \(O = O_1O_2O_T\), and the model \(\lambda\), how do we choose a corresponding state sequence \(Q = q_1q_2...Q_T\) which is optimal in some meaningful sense i.e, best "explains" the observations?}
 Simply put, this problem is about deciphering the most likely hidden states that emitted the visible observation sequences.
 This is done dynamically using the Viterbi algorithm. %%TODO: cite viterbi and explain further if need be
 
\subsubsection{The training problem}
Given the model, \(\lambda\), the training problem raises the following question: \textit{how do we adjust the model parameters \(\lambda\) to maximise the \(P(O|\lambda)\), the probability of the probability of the observation sequence?}
This problem is usually solved by iterative learning algorithms called expectation-maximisation. Examples of this algorithms are Baum-Welch method or any gradient based method. \cite{tuto1989} \cite{tool2001} \cite{biol1998}
\\When using Baum-Welch algorithm, the parameters are initialised by guesses then re-estimated iteratively to find the parameters with maximum likelihood.
This method is vulnerable to local maxima issues. To avoid such cases, it is advice to run it multiple times with different initial values in order to keep the estimation with the highest likelihood value.

\subsubsection{Overfitting, order of markov, robustness: bias-var}

\section{Expection maximisation algorithm} \label{alg:EM}

\section{The Viterbi algorithm} \label{alg:Viterbi}

\section{k-Nearest Neighbour}

\section{Dimensionality reduction}
Dimensionality or dimension reduction is used pattern recognition, machine learning and statistics to find the most compact representation of the dataset by removing redundant and irrelevant information \cite{effi2016}.
It is achieve by extracting principal features, i.e, feature extracting or by selecting the most relevant subset of the initial feature vector, i.e, feature selection, using a supervised or an unsupervised approach.

\subsection{Motivations for dimensionality reduction}
When building a model, the need for dimensionality is supported by several reasons. Some of the important ones are presented in three points.
Firstly, by reducing the feature space's dimension, we can build model with higher quality\cite{newh2016}.
In most classification problems, the feature domains contain variables with very little to no information for the purpose at hand.
Thus, removing these features reduces the complexity of the problem which can in return, increases the model's accuracy.

Secondly, working with hundreds to thousands of features can be diffult to conceptualise and visualise. By using dimensionality reduction, we can better understand the model and present it to others by comprehensive visualisation\cite{newh2016}.

The third reason is about efficiency in terms of computational time and storage. In general, pattern recognition and machine learning algorithms computionally intense. Besides, storage capacity is limited in some engineering applications such as embedded systems. So, solving the problem only with the relevant features can alleviate these two problems. Consequently, the computional speed of the model can increased by using dimensionality reduction \cite{hybr2011}.

Various dimensionality reduction have been developed in literature \cite{sima2013} \cite{effi2016} \cite{hybr2011} \cite{newh2016}, 
the next section will present a handful of the ones used in the present work. 


\subsection{Feature selection: filters and wrappers}
Filters and wrappers are two major categories of feature selection methods \cite{effi2016} \cite{newh2016} \cite{hybr2011} \cite{sima2013}. %%Redraw or crop imagines from original paper
The structure of both methods are illustrated by figure \ref{fig:filt} and figure \ref{fig:wrap} , respectively.
\begin{figure}[ht!]
	\centering
	\includegraphics{Figures/filter}
	\caption{The procedure of filters in dimensionality reduction}{\cite{hybr2011}}
	\label{fig:filt}
\end{figure}
\begin{figure}[ht!]
	\centering
	\includegraphics{Figures/wrapper}
	\caption{The procedure of wrappers in dimensionality reduction}{\cite{hybr2011}}
	\label{fig:wrap}
\end{figure}
They both require a mechanism for generating a subset of the original feature set and a stoppin criterion, which can be a distance measure, a measure of similarity between the features. However, wrappers identify the best feature subset using a learning algorithm usually by fully searching the feature space, whereas, filters use a simple measurement metric based on mutual information, correlation and other distance criteria \cite{hybr2011} \cite{sima2013}.
As a result of the two different approaches, filters are fast and do not guarantee optimal classification accuracy. On the other hand, wrappers give accurate prediction results but are very slow.
Both approaches are often combined to build effective and efficient feature selection methods. One such approach is shown in figure \ref{fig:filt-wrap}. 
\begin{figure}[ht!]
	\centering
	\includegraphics{Figures/filter-wrapper}
	\caption{The procedure of hybrid filter-wrapper in dimensionality reductiony}{\cite{hybr2011}}
	\label{fig:filt-wrap}
\end{figure}
In this approach, the filters are used as a preliminary stage to discard irrelevant features. A feature is deemed irrevant if it cannot discrinate between the different classes or if it contains redundant information \cite{hybr2011}.
The output of the filter stage is a smaller set of relevant features which is fed into the wrapper to find the optimal final subset of features. Thus, the filter stage effectively reduces the search time of the wrapper. 
In classification problems, using the seperability index matrix to determine the classification content or degree of the feature generally results in very good results \cite{sima2013}. 
The next section gives a particular attention to this approach.

\subsubsection{Feature ranking using separability index matrix}
In this section, we dives into a systematic approach to determine the 'classifiability of a feature' as present in \cite{sima2013}.
In their paper, Jeong-Su, Sang Wan Lee and Zeungnam Bien, proposed a new criterion called separability index matrix to identify features that can discriminate between the different classes of a classification problem.
The important concepts of this method are here defined together with their significance.

\begin{enumerate}
	\item \textbf{Separability Degree Matrix (SDM)}\\\\
	\(SDM_k\) is a CxC symmetric matrix of the separability measures between the different classes of classification problem of C distinct classes given a particular feature \(x_k\).
	It is defined in \ref{eq:sdm}
\subitem \[SDM_k = [J(w_i, w_j; {x_k})] \label{eq:sdm}\]	
\subitem where \(J(w_i, w_j; {x_k})\) denotes the separability or the distance value between class \(w_i\)  and class \(w_j\) when the criterion function \(J(\textbf{.})\) such as mahalanobis \cite{maha2013} or eucludian, is applied to the feature \(x_k\) \cite{sima2013}. 
	\(SDM_k\) is symmetric matrix with zero diagonal values, because for a given feature \(x_k\), \(J(w_i, w_j; {x_k}) = J(w_j, w_i; {x_k})\) and \(J(w_i, w_i; {x_k}) = 0\).
	This observation can be exploited  to half the computation required to calculate \(SDM_v = {SDM_1, SDM_2, ..., SDM_N}\), the set of  all \(SDM\) for the feature set of size N.
	 
	\item \textbf{Separability Index Matrix (SIM)}\\\\
	\(SIM_k = {c_{ij}}\) is a CxC matrix of binary values {0, 1}. if \(c_ij=0\), then the classes \(w_i\) and \(w_j\) are not separable by the feature \(x_k\).\\
	\(SIM_k\) is obtained by applying an threshold function to \(SDM_k\).
	For instance, 
	\begin{align*}
		SIM_k(i, j) = 0 \quad if  \quad SDM_k(i, j) < SDM_{avg}(i, j)\\
		SIM_k(i,j) = 1 \quad if  \quad SDM_k(i,j) \geq SDM_{avg}
	\end{align*}
	where \(SDM_{avg}\) is the element wise average of \(SDM_v\).
	\(SIM_k\) is significant because it can be used to systematically determine the irrelevant features, i.e, features whose \(SIM_k = 0\), and/or redundant ones. \cite{sima2013}
	
	
	\item \textbf{Classifiability: \(G(x_k)\)} \label{con:class}\\\\
	Although, conventional distance criterion reveal the separability of the feature distribution, \cite{sima2013}
	 we are often more interested in, how effectively can a particular feature distinguish one class from another. This information is denoted \(G(x_k)\), the ' classifiability of a the feature \(x_k\) '.
	 \(G(x_k)\) is computed by \ref{eq:G}.
	 \begin{align}
	 	G(x_k) = \sum_{i = 1}^{C}\sum_{j = 1}^{C}(SIM_k*WM_k) \label{eq:G}
	 	\\ with \quad WM_k = SIM_k / \sum_{i=1}^{N}SIM_i
	 \end{align}
	 where * and / denote respectively element wise matrix multiplication and division.
\end{enumerate}

This method can be used to effectively and efficiently rank the features in a classification for subset selection or an a preliminary step in a hybrid filter-wrapper feature selection solution.  

\subsection{Feature extraction}
In this work, two distinct feature extraction tools were explored namely, the linear discriminant analysis and the principal component analysis.
Each technique is further explained in the next two subsections.
 
\subsubsection{Linear discriminant analysis}
Linear discrimant analysis maps a K-dimensional dataset, \(O \in \Bbb R^{K}\) into a subspace that maximizes the class between scatter with regards to the class within scatter \cite{prtool}.
In order words, PCA finds the linear transformation - of the dataset into a lower dimensional space - that maximises the between class discrimination \cite{lda2015}
It is a well received dimensionality reduction tool for classification problems by the computer vision and pattern recognition community \cite{lda2015}.

 LDA is a supervised dimensionality reduction method, it does require the dataset to be labeled. On the other hand, PCA, the next feature extraction method is unsupervised method \cite{prtool}.
 
\subsubsection{Principal component analysis}
 Principle component analysis (PCA) is an orthogonal linear transformation of a dataset of T K-dimensional observations, \(O \in \Bbb R^{TxK}\), into a so-called PC-space defined by the principal components (PC) \cite{dime2017}.
 Naturally, the dimension of the new space is smaller or equal to K, the original set's dimensionality. The components are ordered according to the variance of the features of the observation set. The principle components represent the eigenvectors of the covariance matrices of the features and the eigenvalues are measures of the features' variances. Thus, the first component corresponds to the feature with highest variance under the orthogonality constraint \cite{dime2017}.
 In effect, PCA is not transforms the original observation set to a different space but also, ranks the features according to their variability. It can therefore be used as an unsupervised dimensionality reduction technique by simply discarding the features with low variance after finding the principal components.
\include{Design}
\include{Implementation}
\chapter{Results}

\section{Experiement 1: The effect of a CHMM' observation dimensionality on its performance}

\subsection{Aim of the experiement}
The aim of this experiement is to investigate how the number of features impacts the accuracy of a Hidden Markov Model with continuous emission symbols (CHMM), in the abscence of enough training data. Thus, the hypothesis under investigation is:\\
\textbf{\textit{In the absence of enough training data, a CHMM with observations of high dimensionality performs poorly.}}

\subsection{Experiment apparatus}
To perform this expereiment, the following materials are required:
\begin{itemize}
	\item \(\lambda\), a continuous Hidden Markov Model specified by \(\lambda = (A, \beta_{jm}, \mu_{jm}, \Sigma_{jm}, \pi)\).
	\item At least two sample data sets training the model and testing it.
	\item A criterion to rank and select subsets of features.
	\item A measure to evaluate the performance of the CHMM model.
	\item Finally, a way to visualise the results of the experiments
\end{itemize}

\subsection{Experiment procedure}
The expreriment was performed with the steps listed below:
\begin{enumerate}
	\item Step 0 - Preliminary data pre-processing: This step consisted in the data pre-processing as described in %%TODO: cross-ref data pre-processing section
	\item Step 1 - Partitioning data into training and test sets: Here, the dataset was randomly sampled into training and test sets. The training set was relatively small, it was a sequence 539 observations. 
	\item Step 2 -  Feature ranking: The features were sorted in a descending order based on their ability to discriminate the different states of the CHMM. The separability index method described in was used for this purpose. %%TODO: quote et cross-ref SIM
	\item Step 3 - Data subset selection: Select the optimal feature subset, starting with 1 dimension. 
	\item Step 4 - Model building and training: The CHMM model, \(\lambda\) was built and trained using with training dataset using the optimal feature subset.
	\item Step 5 - Model testing and evaluation: The model was tested with the test dataset. The test consisted in decoding the most likely state sequence given a previously unseen sequence of observations. This path prediction was evaluated based on the evaluation criterion presented in %%TODO: cross-reference evaluation criterion.
	\item Step 6 - Iteration: Step 3 through step 5 were repeated while varying the feature subset size until the maximum size, which is 18 in this case. In each iteration, the prediction accuracy was stored in an array for visualisation.
	\item The different accuracies were finally ploted as a function of the observation dimensionality. Moreover, the observations were grouped based on the corresponding hidden state sequence and scattered in a 2-dimensional principal component space. This is to compare the decoded states against the ground-truth.
\end{enumerate}

\subsection{Experiment results}
The results of the experiments are presented in figure \ref{fig:dim-acc}, \ref{fig:dim-log}, \ref{fig:gt-5dim}, \ref{fig:es-5dim}, \ref{fig:gt-18dim}, \ref{fig:es-18dim}.\\
\ref{fig:dim-acc} and \ref{fig:dim-log} show how the hidden state decoding performance and the log-likelihood estimated by the CHMM model, train with just 539 observations, varies as the observation dimensionality increases.\\
\ref{fig:gt-5dim} and \ref{fig:es-5dim}, \ref{fig:gt-18dim} and \ref{fig:es-18dim} illustrate how the estimated state sequences compare to reality for observation sequences of 5-dimensions and 18-dimensions.    
5 and 18 dimensions were presented because they are the two extremes in terms of accuracy. Other dimensions may be found in the appendices, %%TODO: ref appendices
\begin{figure}[ht!]
	\includegraphics[width=\textwidth,height=\textheight,keepaspectratio]{Figures/dimensionality-effect-acc}
	\caption{The effect of CHMM's observation dimensionality the state sequence decoding accuracy}
	\label{fig:dim-acc}
\end{figure}

\begin{figure}[ht!]
	\includegraphics{Figures/dimensionality-effect-log}
	\caption{The effect of CHMM's observation dimensionality the log-likelihood}
	\label{fig:dim-log}
\end{figure}


\begin{figure}[ht!]
	\includegraphics{Figures/ground-truth-scatter-with-5-features}
	\caption{Scatter plot of  5-dimensional observations grouped per state based on ground-truth state sequence}
	\label{fig:gt-5dim}
\end{figure}
\begin{figure}[ht!]
	\includegraphics{Figures/estimation-scatter-with-5-features}
	\caption{Scatter plot of 5-dimensional observations grouped per state based on estimated state sequence}
	\label{fig:es-5dim}
\end{figure}

\begin{figure}[ht!]
	\includegraphics{Figures/ground-truth-scatter-with-18-features}
	\caption{Scatter plot of 18-dimensional observations grouped per state based on ground-truth state sequence}
	\label{fig:gt-18dim}
\end{figure}

\begin{figure}[ht!]
	\includegraphics{Figures/estimation-scatter-with-18-features}
	\caption{Scatter plot of 18-dimensional observations grouped per state based on estimated state sequence}
	\label{fig:es-18dim}
\end{figure}

\subsection{Analysis of results}

\subsection{Conclusions and recommendations of the experiment}

\section{Experiement 2: The effect of dimensionality reduction on CHMM's performance}

\subsection{Aim of the experiement}
The aim of this experiment is to investigate the effect of dimensionality reduction on the performance of a continuous Hidden Markov Model (CHMM).
The hypothesis under investigation is therefore the following:
\textbf{\textit{Dimension reduction can cause an increase in a CHMM model performance when there is not enough training data.}}
Thus, the performance of the CHMM without and with various dimensionality reduction are compared to test the hypothesis.

\subsection{Experiment apparatus}
The assets needed to perform the experiment are listed below.
\begin{itemize}
	\item \(\lambda\), a continuous Hidden Markov Model specified by \(\lambda = (A, \beta_{jm}, \mu_{jm}, \Sigma_{jm}, \pi)\).
	\item Dimensionality reduction methods. Two wrapper and two filter methods were considered. The wrapper methods were Principle Component Analysis (PCA), Linear Discriminant Analysis (LDA). The two filters methods were feature ranking with similarity index %%TODO cross-ref
	and a combination of forward feature selection and similarity index.
	\item A Performance metric. The metric used were the hidden state decoding accuracy and the log-likelihood of the EM algorithm.
\end{itemize}

\subsection{Experiment procedure}
The experiment was performed as follows.
First, the dataset was partition into two different set for training and testing using random sampling.
Using the same traing dataset five different models were built and trained, \(\Lambda = (\lambda_{No Reduction}, \lambda_{PCA}, \lambda_{LDA}, \lambda_{SI}, \lambda_{SI-forward})\).
Then the different models were tested with the same test dataset and the prediction accuracy as well as the EM algorithm loglikehood were recorded.
The training and the testing were repeated while varying the proportion of training set used from \(10\%\) to \(90\%\) of the total dataset.
The prediction accuracies and the EM algorithm likelihoods were finally plotted as a function of the training data size for each model.
These findings are presented in the figures \ref{fig:size-acc} and \ref{fig:size-log}. 

\subsection{Experiment results}
Firstly, figure \ref{fig:size-acc} how the performances of the five CHMMs compare against each other as the training data size increases.
Secondly, the loglikelihoods presented in \ref{fig:size-log} show how effectively can each model recognise an observation sequence generated by the underlying mechanism.

\begin{figure}[ht!]
	\includegraphics{Figures/size-acc-2}
	\caption{The effect of training datasize on the prediction accuracy}
	\label{fig:size-acc}
\end{figure}

\begin{figure}[ht!]
	\includegraphics{Figures/size-log-2}
	\caption{The effect of training datasize on the log likelihood}
	\label{fig:size-log}
\end{figure}

\begin{figure}[ht!]
	\centering
	\includegraphics{Figures/bias-var}
	\caption{Bias-Variance tradeoff analysis}
	\label{fig:bias-var}
\end{figure}


\begin{figure}[ht!]
	%\includegraphics{Figures/dimensionality-effect-log}
	%\caption{The effect of CHMM's observation dimensionality the log-likelihood}
	%\label{fig:dim-log}
\end{figure}


\begin{figure}[ht!]
	%\includegraphics{Figures/ground-truth-scatter-with-5-features}
	%\caption{Scatter plot of  5-dimensional observations grouped per state based on ground-truth state sequence}
	%\label{fig:gt-5dim}
\end{figure}
\begin{figure}[ht!]
	%\includegraphics{Figures/estimation-scatter-with-5-features}
	%\caption{Scatter plot of 5-dimensional observations grouped per state based on estimated state sequence}
	%\label{fig:es-5dim}
\end{figure}

\begin{figure}[ht!]
	%\includegraphics{Figures/ground-truth-scatter-with-18-features}
	%\caption{Scatter plot of 18-dimensional observations grouped per state based on ground-truth state sequence}
	%\label{fig:gt-18dim}
\end{figure}

\begin{figure}[ht!]
	%\includegraphics{Figures/estimation-scatter-with-18-features}
	%\caption{Scatter plot of 18-dimensional observations grouped per state based on estimated state sequence}
	%\label{fig:es-18dim}
\end{figure}

\subsection{Analysis of results}

\subsection{Conclusions and recommendations of the experiment}
\chapter{Discussion}

Here is what the results mean and how they tie to existing literature...

Discuss the relevance of your results and how they fit into the theoretical work you described in your
literature review.

\chapter{General conclusions}
The present report consolidates the work done on gait sequence modelling and estimation using Hidden Markov Models (HMMs). 
We went from the continuous IMU measurements of a dog moving, to gait state identification using HMM with Gaussian mixture density functions.\\
Given limitation imposed by the small size of the dataset, dimensionality reduction techniques were explored. In addition, data aggregation and mirroring techniques were employed to increase the data. After successfully implementing the various models, experiments were designed and performed to test the hypotheses formulated in the introduction. These hypotheses are repeated here as a reminder.
\begin{enumerate}
	\item HMMs can successfully model gait sequence dynamics using IMU data, in the absence of huge training data.
	\item Data aggregation and mirroring techniques can be used to overcome training data size limitations.
	\item Dimensionality reduction can increase the performance of an HMM, in the absence of a large training set.
	\item HMMs can be used to successfully perform gait activity recognition.
\end{enumerate}
In light of the above results and the discussions, these hypotheses can be confidently accepted. \\
Indeed, continuous HMMs with Gaussian mixture can be used to build robust gait sequence estimators from IMI measurements with up to 95\% precision.
This accuracy is obtainable by aggregating the measurements of the different gait actions and/or by using considering the reverse gait sequence as a valid sequence.\\
Furthermore, when dealing with a small training data, dimensionality reduction can increase the model's performance by up to 78\%. However, with a large dataset, it might not be necessary to perform any reduction with just 18-dimensions. This last conclusion should be taken with some degree of reservation.
Finally, the last experiment showed that IMU based HMM can be used for gait type classification.\\
The objectives of this project can, therefore, be considered met. In the next section recommendations are made for possible improvements and future works.
\chapter{Recommendations}
In this chapter are presented the recommendations that follow from the completion of the project. These recommendations can serve as an extension of the project for improvement and future work.
\begin{itemize}
	\item \textit{Online real-time testing:} Even though the model was tested with unseen data, an online real-time implementation and testing on a moving dog would testify of the reliability of the model built.
	
	\item \textit{Gather more data from different dogs:} The HMM model was formulated and the experiments were performed based on measurements from a single subject. More data should be gathered from multiple dogs to test the present and subsequently tune the model's parameters based on the outcome of the tests.
	
	\item \textit{Incorporate more domain knowledge:} Given the limited amount of time, domain-specific knowledge was not thoroughly explored. Making more use of available knowledge on quadruped movement may improve the model's complexity and performance. 
	
	\item \textit{Compare the HMM models to other algorithms:} The performance of the constructed models should be compared to other classification and pattern recognition methods.
	
	\item \textit{Combine the front and back HMMs:} In this design, separate models were built for the front and the back legs of the dog. Although this made the algorithm applicable to bipedal gait estimation, it would be useful to combined the two models and evaluate the performance of the holistic 16-states HMM.
	
	\item \textit{Investigate more dimensionality reduction techniques:} The feature extraction methods performed did not perform very well, more dimensionality reduction methods such as Fast Fourier Transform (FFT) \cite{towa2009}, time-frequency analysis \cite{ches2012} should be investigated.
\end{itemize}



\begin{thebibliography}{5}
	
\bibitem{tuto1989} L. R. Rabiner, "A tutorial on hidden Markov models and selected applications in speech recognition", \emph{Proceedings of the IEEE}, vol. 77, pp. 257-286, 1989, ISSN 0018-9456.

\bibitem{tool2001} Huseyin M. Ertunc and Kenneth A. Loparo and Hasan Ocak, "Tool wear condition monitoring in drilling operations using hidden Markov models (HMMs)", \emph{International Journal of Machine Tools and Manufacture}, vol. 41, no. 9, pp. 1363-1384, 2001, ISSN 0018-9456, DOI https://doi.org/10.1016/S0890-6955(00)00112-7. [Online]. Available:  http://www.sciencedirect.com/science/article/pii/S0890695500001127

\bibitem{cont2013} Ghazaleh Panahandeh, Nasser Mohammadiha, Arne Leijon, Peter Händel, "Continuous Hidden Markov Model for Pedestrian Activity Classification and Gait Analysis", \emph{Instrumentation and Measurement IEEE Transactions on}, vol. 62, pp. 1073-1083, 2013, ISSN 0890-6955, DOI 10.1109/5.18626

\bibitem{towa2009} Guangyi Shi and Yuexian Zou and Yufeng Jin and Xiaole Cui and W. J. Li, "Towards HMM based human motion recognition using MEMS inertial sensors", \emph{2008 IEEE International Conference on Robotics and Biomimetics}, pp. 1762-1766, 2009, DOI 10.1109/ROBIO.2009.4913268.

\bibitem{ches2012} G. Panahandeh and N. Mohammadiha and A. Leijon and P. Händel, "Chest-mounted inertial measurement unit for pedestrian motion classification using continuous hidden Markov model", \emph{2012 IEEE International Instrumentation and Measurement Technology Conference Proceedings}, pp. 991-995, 2012, ISSN 1091-5281, DOI 10.1109/I2MTC.2012.6229380.

\bibitem{effi2016} Mingxia Liu and Daoqiang Zhang, "Feature selection with effective distance", \emph{Neurocomputing}, vol. 215, no. Supplement C, pp. 100-109, 2016, note SI: Stereo Data, ISSN 0925-2312, DOI https://doi.org/10.1016/j.neucom.2015.07.155. [Online]. Available: http://www.sciencedirect.com/science/article/pii/S0925231216306336. 

\bibitem{twop2008} Ali Mahvash and Annie Ross, "Two-phase flow pattern identification using continuous hidden Markov model", \emph{International Journal of Multiphase Flow}, vol. 34, no. 3, pp. 303-311, 2008, ISSN 0301-9322, DOI https://doi.org/10.1016/j.ijmultiphaseflow.2007.08.006. [Online]. Available: http://www.sciencedirect.com/science/article/pii/S0301932207001462. 

\bibitem{biol1998} R. Durbin and S. Eddy and A. Krogh and G. Mitchison, Markov chains and hidden Markov models", in \emph{Biological sequence analysis}, UK: Cambridge University Press, 1998, pp. 46-79.

\bibitem{newh2016} Saúl Solorio-Fernández and J. Ariel Carrasco-Ochoa and José Fco. Martínez-Trinidad, "A new hybrid filter–wrapper feature selection method for clustering based on ranking", \emph{Neurocomputing}, vol. 214, no. Supplement C, pp. 866-880, 2016, ISSN 0925-2312, DOI https://doi.org/10.1016/j.neucom.2016.07.026. [Online]. Available: http://www.sciencedirect.com/science/article/pii/S0925231216307718. 

\bibitem{hybr2011} Hui-Huang Hsu and Cheng-Wei Hsieh and Ming-Da Lu, "Hybrid feature selection by combining filters and wrappers", \emph{Expert Systems with Applications}, vol. 38, no. 7, pp. 8144-8150, 2011, ISSN 0957-4174, DOI https://doi.org/10.1016/j.eswa.2010.12.156. [Online]. Available: http://www.sciencedirect.com/science/article/pii/S0957417410015198. 

\bibitem{sima2013} Jeong-Su Han and Sang Wan Lee and Zeungnam Bien, "Feature subset selection using separability index matrix", \emph{Information Sciences}, vol. 223, no. Supplement C, pp. 102-118, 2013, ISSN 0020-0255, DOI https://doi.org/10.1016/j.ins.2012.09.042. [Online]. Available: http://www.sciencedirect.com/science/article/pii/S0020025512006354. 

\bibitem{maha2013} Zhao Yongli and Zhang Yungui and Tong Weiming and Chen Hongzhi, "An improved feature selection algorithm based on MAHALANOBIS distance for Network Intrusion Detection", \emph{PROCEEDINGS OF 2013 International Conference on Sensor Network Security Technology and Privacy Communication System}, pp. 69-73, 2013, DOI 10.1109/SNS-PCS.2013.6553837.

\bibitem{prtool} Dimension Reduction - Pattern Recognition Tools.[Online]. Available: http://37steps.com/prtools/examples/dimension-reduction/

\bibitem{lda2015} Youness Aliyari Ghassabeh and Frank Rudzicz and Hamid Abrishami Moghaddam, "Fast incremental LDA feature extraction", \emph{Pattern Recognition}, vol. 48, no. 6, pp. 1999-2012, 2015, ISSN 0031-3203, DOI https://doi.org/10.1016/j.patcog.2014.12.012. [Online]. Available: http://www.sciencedirect.com/science/article/pii/S0031320314005214.

\bibitem{dime2017} Miguel Simão and Pedro Neto and Olivier Gibaru, "Using data dimensionality reduction for recognition of incomplete dynamic gestures", \emph{Pattern Recognition Letters}, vol. 99, no. Supplement C, pp. 32-38, 2017, ISSN 0167-8655, DOI https://doi.org/10.1016/j.patrec.2017.01.003. [Online]. Available: http://www.sciencedirect.com/science/article/pii/S016786551730003X.

\bibitem{pendu2004} Timothy M. Griffin and Russell P. Main and Claire T. Farley, "Biomechanics of quadrupedal walking: how do four-legged animals achieve
inverted pendulum-like movements?", \emph{The Journal of Experimental Biology}, vol. 207, pp. 3545-3558, 2004, DOI 10.1242/jeb.01177.

\bibitem{mat-feat} Feature Selection - MathWorks.[Online]. Available: https://www.mathworks.com/help/stats/feature-selection.html

\bibitem{knnc} K-Nearest Neighbor Classifier - Pattern Recognition Tools.[Online]. Available: http://www.37steps.com/prhtml/prtools/knnc.html

\bibitem{hmmestimate} hmmestimate - MathWorks .[Online]. Available: https://www.mathworks.com/help/stats/hmmestimate.html

\bibitem{gmm2000} McLachlan, G., and D. Peel, Finite Mixture Models. Hoboken, NJ: \emph{John Wiley \& Sons}, Inc., 2000.

\bibitem{fitgmm} fitgmdist - MathWorks .[Online]. Available:  https://www.mathworks.com/help/stats/fitgmdist.html

\bibitem{AIC_ref} Ljung, L. System Identification: Theory for the User, Upper Saddle River, NJ, Prentice-Hall PTR, 1999

\bibitem{AIC} Akaike's Information Criterion for estimated model - MATLAB aic .[Online]. Available:  https://www.mathworks.com/help/ident/ref/aic.html

\bibitem{BIC} Bayes Information Criterion - MATLAB - MathWorks .[Online]. Available: https://www.mathworks.com/help/stats/gmdistribution.bic.html

\bibitem{matlab-hmm} Qiuqiang Kong, "matlab-hmm" .[Online]. Available: https://github.com/qiuqiangkong/matlab-hmm

\bibitem{featsel} Trainable mapping for forward feature selection - Pattern Recognition Tools. [Online] . http://www.37steps.com/prhtml/prtools/featself.html

\bibitem{maha} Mahalanobis distance - Pattern Recognition Tools. [Online] http://www.37steps.com/prhtml/prtools/distmaha.html

\bibitem{pcam} Principal component analysis (PCA or MCA on overall covariance matrix) - Pattern Recognition Tools. [Online] http://www.37steps.com/prhtml/prtools/pcam.html

\bibitem{ldam} Optimal discrimination linear mapping (Fisher mapping, LDA) - Pattern Recognition Tools. [Online]. http://www.37steps.com/prhtml/prtools/fisherm.html

\bibitem{cvpartition} Create cross validation partition for data - MATLAB - MathWorks .[Online]. https://www.mathworks.com/help/stats/cvpartition.html

\bibitem{gendat} Random sampling of datasets for training and testing - Pattern Recognition Tools. [Online]. http://www.37steps.com/prhtml/prtools/gendat.html

\bibitem{prtool-home} PRTools - A Matlab toolbox for pattern recognition .[Online]. http://prtools.org/

\bibitem{matlab-stats} Statistics and Machine Learning Toolbox - MATLAB - MathWorks .[Online]. https://www.mathworks.com/products/statistics.html
\end{thebibliography}
\appendix
\chapter{Additional Files and Schematics}

\section{Implementation and experiment code} \label{apdix-repo}
The implementation can be accessed from this github repository: \href{https://github.com/h-kouame/gait-sequence-modelling-and-estimation-v2}{https://github.com/h-kouame/gait-sequence-modelling-and-estimation-v2}

\section{Additional results} \label{apdix-results}

\begin{figure}[ht!]
	\includegraphics{Figures/ground-truth-scatter-with-1-features}
	\caption{Scatter plot of  1-dimensional observations grouped according the ground-truth state sequence}
	\label{fig:gt-1dim}
\end{figure}
\begin{figure}[ht!]
	\includegraphics{Figures/estimation-scatter-with-1-features}
	\caption{Scatter plot of 1-dimensional observations grouped according to the estimated state sequence}
	\label{fig:es-1dim}
\end{figure}

\begin{figure}[ht!]
	\includegraphics{Figures/ground-truth-scatter-with-2-features}
	\caption{Scatter plot of 2-dimensional observations grouped according to the ground-truth state sequence}
	\label{fig:gt-2dim}
\end{figure}

\begin{figure}[ht!]
	\includegraphics{Figures/estimation-scatter-with-2-features}
	\caption{Scatter plot of 2-dimensional observations grouped according to the estimated state sequence}
	\label{fig:es-2dim}
\end{figure}

\begin{figure}[ht!]
	\includegraphics{Figures/ground-truth-scatter-with-3-features}
	\caption{Scatter plot of  3-dimensional observations grouped according the ground-truth state sequence}
	\label{fig:gt-3dim}
\end{figure}
\begin{figure}[ht!]
	\includegraphics{Figures/estimation-scatter-with-3-features}
	\caption{Scatter plot of 3-dimensional observations grouped according to the estimated state sequence}
	\label{fig:es-3dim}
\end{figure}

\begin{figure}[ht!]
	\includegraphics{Figures/ground-truth-scatter-with-4-features}
	\caption{Scatter plot of 4-dimensional observations grouped according to the ground-truth state sequence}
	\label{fig:gt-4dim}
\end{figure}

\begin{figure}[ht!]
	\includegraphics{Figures/estimation-scatter-with-4-features}
	\caption{Scatter plot of 4-dimensional observations grouped according to the estimated state sequence}
	\label{fig:es-4dim}
\end{figure}

\begin{figure}[ht!]
	\includegraphics{Figures/ground-truth-scatter-with-6-features}
	\caption{Scatter plot of  6-dimensional observations grouped according the ground-truth state sequence}
	\label{fig:gt-6dim}
\end{figure}
\begin{figure}[ht!]
	\includegraphics{Figures/estimation-scatter-with-6-features}
	\caption{Scatter plot of 6-dimensional observations grouped according to the estimated state sequence}
	\label{fig:es-6dim}
\end{figure}

\begin{figure}[ht!]
	\includegraphics{Figures/ground-truth-scatter-with-7-features}
	\caption{Scatter plot of 7-dimensional observations grouped according to the ground-truth state sequence}
	\label{fig:gt-7dim}
\end{figure}

\begin{figure}[ht!]
	\includegraphics{Figures/estimation-scatter-with-7-features}
	\caption{Scatter plot of 7-dimensional observations grouped according to the estimated state sequence}
	\label{fig:es-7dim}
\end{figure}

\begin{figure}[ht!]
	\includegraphics{Figures/ground-truth-scatter-with-8-features}
	\caption{Scatter plot of  8-dimensional observations grouped according the ground-truth state sequence}
	\label{fig:gt-8dim}
\end{figure}
\begin{figure}[ht!]
	\includegraphics{Figures/estimation-scatter-with-8-features}
	\caption{Scatter plot of 8-dimensional observations grouped according to the estimated state sequence}
	\label{fig:es-8dim}
\end{figure}

\begin{figure}[ht!]
	\includegraphics{Figures/ground-truth-scatter-with-9-features}
	\caption{Scatter plot of 9-dimensional observations grouped according to the ground-truth state sequence}
	\label{fig:gt-9dim}
\end{figure}

\begin{figure}[ht!]
	\includegraphics{Figures/estimation-scatter-with-9-features}
	\caption{Scatter plot of 9-dimensional observations grouped according to the estimated state sequence}
	\label{fig:es-9dim}
\end{figure}
	
	\begin{figure}[ht!]
		\includegraphics{Figures/ground-truth-scatter-with-10-features}
		\caption{Scatter plot of 10-dimensional observations grouped according to the ground-truth state sequence}
		\label{fig:gt-10dim}
	\end{figure}
	
	\begin{figure}[ht!]
		\includegraphics{Figures/estimation-scatter-with-10-features}
		\caption{Scatter plot of 10-dimensional observations grouped according to the estimated state sequence}
		\label{fig:es-10dim}
	\end{figure}
	
	\begin{figure}[ht!]
		\includegraphics{Figures/ground-truth-scatter-with-11-features}
		\caption{Scatter plot of  11-dimensional observations grouped according the ground-truth state sequence}
		\label{fig:gt-11dim}
	\end{figure}
	\begin{figure}[ht!]
		\includegraphics{Figures/estimation-scatter-with-11-features}
		\caption{Scatter plot of 11-dimensional observations grouped according to the estimated state sequence}
		\label{fig:es-11dim}
	\end{figure}
	
	\begin{figure}[ht!]
		\includegraphics{Figures/ground-truth-scatter-with-12-features}
		\caption{Scatter plot of 12-dimensional observations grouped according to the ground-truth state sequence}
		\label{fig:gt-12dim}
	\end{figure}
	
	\begin{figure}[ht!]
		\includegraphics{Figures/estimation-scatter-with-12-features}
		\caption{Scatter plot of 12-dimensional observations grouped according to the estimated state sequence}
		\label{fig:es-12dim}
	\end{figure}
	
	\begin{figure}[ht!]
		\includegraphics{Figures/ground-truth-scatter-with-13-features}
		\caption{Scatter plot of  13-dimensional observations grouped according the ground-truth state sequence}
		\label{fig:gt-13dim}
	\end{figure}
	\begin{figure}[ht!]
		\includegraphics{Figures/estimation-scatter-with-13-features}
		\caption{Scatter plot of 13-dimensional observations grouped according to the estimated state sequence}
		\label{fig:es-13dim}
	\end{figure}
	
	\begin{figure}[ht!]
		\includegraphics{Figures/ground-truth-scatter-with-14-features}
		\caption{Scatter plot of 14-dimensional observations grouped according to the ground-truth state sequence}
		\label{fig:gt-14dim}
	\end{figure}
	
	\begin{figure}[ht!]
		\includegraphics{Figures/estimation-scatter-with-14-features}
		\caption{Scatter plot of 14-dimensional observations grouped according to the estimated state sequence}
		\label{fig:es-14dim}
	\end{figure}
	
	\begin{figure}[ht!]
		\includegraphics{Figures/ground-truth-scatter-with-15-features}
		\caption{Scatter plot of  15-dimensional observations grouped according the ground-truth state sequence}
		\label{fig:gt-15dim}
	\end{figure}
	\begin{figure}[ht!]
		\includegraphics{Figures/estimation-scatter-with-15-features}
		\caption{Scatter plot of 15-dimensional observations grouped according to the estimated state sequence}
		\label{fig:es-15dim}
	\end{figure}
	
	\begin{figure}[ht!]
		\includegraphics{Figures/ground-truth-scatter-with-16-features}
		\caption{Scatter plot of 16-dimensional observations grouped according to the ground-truth state sequence}
		\label{fig:gt-16dim}
	\end{figure}
	
	\begin{figure}[ht!]
		\includegraphics{Figures/estimation-scatter-with-16-features}
		\caption{Scatter plot of 16-dimensional observations grouped according to the estimated state sequence}
		\label{fig:es-16dim}
	\end{figure}

	\begin{figure}[ht!]
	\includegraphics{Figures/ground-truth-scatter-with-17-features}
	\caption{Scatter plot of 17-dimensional observations grouped according to the ground-truth state sequence}
	\label{fig:gt-17dim}
	\end{figure}

	\begin{figure}[ht!]
		\includegraphics{Figures/estimation-scatter-with-17-features}
		\caption{Scatter plot of 17-dimensional observations grouped according to the estimated state sequence}
		\label{fig:es-17dim}
	\end{figure}

\chapter{Addenda}

\section{Ethics Forms}
\begin{figure}[ht!]
	\includegraphics[width=1.5\textwidth, angle=90]{Figures/ethics-form}
%	\caption{The effect of training datasize on the log likelihood}
%	\label{fig:size-log}
\end{figure}
}
\end{document}
